% The three time words this node deals in, drawn to the same scale, with the
% range each one buys.
% Note for anyone copying this figure: \framefield always draws at y = 0, so
% each row lives in its own shifted scope rather than carrying a y coordinate.
\begin{tikzpicture}[font=\sffamily\small]

% ================================================================ the node's own
\begin{scope}
  \node[signame] at (-2mm,3mm) {this node};
  \framefield{0}{34}{fid}{upper 32 bits\\{\tiny software, from the update interrupt}}
  \framefield{34}{34}{fdata}{lower 32 bits\\{\tiny TIM2->CNT, one tick is 1 us}}
  % \framebits places its label at 6 mm, which is inside a box this tall, so
  % the bit ranges are written directly above each field instead.
  \node[fbyte] at (17mm,12.5mm) {63 .. 32};
  \node[fbyte] at (51mm,12.5mm) {31 .. 0};
  \node[chlabel, anchor=west, text width=52mm] at (72mm,3mm)
    {64 bits of microseconds: 584,542 years. Monotonic, never stepped, never
     reset. The only word this node computes with.};
\end{scope}

% ================================================================ the wrap zoo
\begin{scope}[yshift=-26mm]
  \node[signame] at (-2mm,4mm) {what wraps, and when};
  \draw[busline] (0mm,4mm) -- (68mm,4mm);
  \foreach \x/\lab/\t in {4/{16 bit at 1 us}/{65.5 ms},
                          22/{16 bit bus stamp}/{seconds},
                          42/{32 bit at 3.57 ns}/{15.3 s},
                          62/{32 bit at 1 us}/{71.6 min}}{
    \draw[busstub] (\x mm,4mm) -- (\x mm,0mm);
    \node[font=\sffamily\tiny, anchor=north] at (\x mm,0mm) {\lab};
    \node[font=\sffamily\tiny, anchor=north, text=red!60!black] at (\x mm,-3.5mm) {\t};
  }
  \node[chlabel, anchor=west, text width=52mm] at (72mm,2mm)
    {Every one of these is extended in software, or it is a defect waiting for
     an inconvenient moment.};
\end{scope}

% ================================================================ the profile's word
\begin{scope}[yshift=-52mm]
  \node[signame] at (-2mm,3mm) {field bus profile};
  \framefield{0}{30}{fdata}{milliseconds after midnight}
  \framefield{30}{8}{fstuff}{void}
  \framefield{38}{30}{fcrc}{days since 1 January 1984}
  \node[fbyte] at (15mm,9mm) {27 .. 0};
  \node[fbyte] at (34mm,9mm) {31 .. 28};
  \node[fbyte] at (53mm,9mm) {47 .. 32};
  \node[chlabel, anchor=west, text width=52mm] at (72mm,3mm)
    {48 bits, and it is a wall clock: a date, not an elapsed time. The epoch is
     Sunday 1 January 1984.};
\end{scope}

% ================================================================ the third one
\begin{scope}[yshift=-72mm]
  \node[signame] at (-2mm,3mm) {profile, high resolution};
  \framefield{0}{68}{fdata}{microseconds, unsigned, 32 bits}
  \node[fbyte] at (34mm,7mm) {31 .. 0};
  \node[chlabel, anchor=west, text width=52mm] at (72mm,3mm)
    {Wraps about every 71.6 minutes. Handle the wrap, or do not put it in a
     frame at all.};
\end{scope}

% ================================================================ the reading
\node[umlnote, text width=60mm, anchor=north west] at (0mm,-84mm)
  {Three different words for three different jobs, and the top one is the only
   one this node computes with. The two profile words exist because chapter 11
   has to put a timestamp in a frame that another vendor's tool can read, and
   the mapping between the node's own 64-bit microseconds and those two is
   arithmetic that belongs in one file with tests, not scattered through the bus
   code.};
\node[umlnote, text width=52mm, anchor=north west] at (66mm,-84mm)
  {The second row is the reason this chapter comes before the bus chapters.
   Every hardware counter here wraps inside an hour, two of them inside a
   minute, and a joint node runs for weeks.};
\end{tikzpicture}
